\begin{claim}{C.71}{Proved}{Actual disk comparisons for two-dimensional operators and a cyclic compactification}
Let $Q$ be a closed locally compact subgroup of
$\operatorname{Homeo}(\mathbb D)$, with its actual compact-open
and inverse topology, acting freely and properly on $\mathbb D$.
When $\dim Q=2$, put
\[
 E_Q=\{h\in\mathcal D_\partial:
                    h|_{\mathbb D}\in
                    C_{\operatorname{Homeo}(\mathbb D)}(Q)\}.
                                                               \tag{C.71.1}
\]
Then every based sphere family $S^n\to E_Q$, $n\geq1$, has a
based contraction in $E_Q$. In particular every based loop
fills there. No extension of the $Q$-action to the boundary,
and no closed-disk conjugacy for that action, is assumed.
If $A\leq\operatorname{Homeo}(\overline{\mathbb D})$ has
actual restricted closure $Q$, then
$E_Q=C_{\mathcal D_\partial}(A)$ by continuity on the
interior and density in the closed disk.

There is also an explicit cyclic action on the actual closed
disk whose interior image is closed, discrete, free and proper,
and which admits no closed-disk
conjugacy to a torsion-free cyclic Fuchsian action. For its
specified invariant partial compactification, the analogue of
the restriction isomorphism (C.69.2) is not surjective. The
proof gives the actual maps witnessing both assertions.
Nevertheless its boundary-fixed centralizer is contractible:
it is identified topologically with an explicitly described
group of continuous families of increasing real lifts, and
linear interpolation supplies the contraction. Its interior
orbits are precisely the displayed vertical lines, so this
example does not supply an open transitive orbit for that
boundary-fixed centralizer.
\end{claim}

\begin{proof}
\emph{The actual two-dimensional transformation group.}
The Lie conclusion and principal orbit comparison of C.33
apply to $Q$. The identity component has open orbits because
the free orbit dimension is two. All these orbits are open;
connectedness of $\mathbb D$ leaves just one. Given $q\in Q$,
transitivity supplies $q_0\in Q^0$ with $q_0x=qx$ at one
specified $x$. Freeness gives $q=q_0$. Thus $Q=Q^0$, and
evaluation
\[
                 \epsilon_x:Q\longrightarrow\mathbb D,
                       \qquad q\longmapsto qx
                                                               \tag{C.71.2}
\]
is a homeomorphism. The inverse is the principal orbit
coordinate from C.33; alternatively properness makes the
continuous orbit bijection a homeomorphism. In particular
$Q$ is simply connected.

Its two-dimensional Lie algebra is either abelian or has
a basis $X,Y$ with $[X,Y]=Y$: in the nonabelian case the
derived algebra is a line and rescaling a complement gives
that bracket. The simply connected groups are respectively
$\mathbb R^2$ and
\[
 (a,b)(a',b')=(a+a',b+e^a b'),
                   \qquad(a,b)\in\mathbb R^2 .
                                                               \tag{C.71.3}
\]
This uses the actual Lie-group isomorphisms integrating those
bases. Their inverses are the integrations in the opposite
direction; the uniqueness for simply connected groups gives
both inverse identities. In each case exponential is a
global diffeomorphism. In (C.71.3) its explicit formula is
$\exp(a,v)=(a,v(e^a-1)/a)$, with the quotient at $a=0$
defined to be $1$.

Through (C.71.2), a map commuting with all left translations
is precisely a right translation. Indeed its value $c$ at
the identity forces its value at every $q$ to be $qc$.
Right translations and their inverses depend continuously
on $c$ in the compact-open topology, and evaluation gives
the inverse comparison. Thus the actual interior centralizer
is topologically the opposite group
\[
 N=C_{\operatorname{Homeo}(\mathbb D)}(Q)\cong Q^{\rm op}.
                                                               \tag{C.71.4}
\]
Inversion is an isomorphism $Q^{\rm op}\to Q$, so $N$ has
the same two specified simply connected exponential models.

\emph{Compact words retain the given disk boundary.}
Let $a:S^n\to E_Q$ be based, with $n\geq1$, and let
$g_z=a(z)|_{\mathbb D}\in N$. Define the interior maps
\[
               R_{r,z}=\exp_N(r\log_N g_z),
                     \qquad 0\leq r\leq1 .
                                                               \tag{C.71.5}
\]
They are continuous in the actual interior homeomorphism
topology. They equal the identity for $r=0$ and at the
sphere basepoint, and equal $g_z$ for $r=1$.

Put
\[
 P=a(S^n)\cup a(S^n)^{-1}\subset\mathcal D_\partial .
                                                               \tag{C.71.6}
\]
This is compact and contains the identity. For each $z$,
choose a path in $S^n$ from its basepoint to $z$.
Applying the exact two- or three-letter formula (C.37.2)
to its restricted path in the simply connected group $N$
places $R_{r,z}$ in the restriction of $P^3$.
No continuous choice of these sphere paths is used.
The resulting restriction map
\[
                  P^3\longrightarrow
                         \operatorname{Homeo}(\mathbb D)
                                                               \tag{C.71.7}
\]
is continuous and injective: two closed-disk maps agreeing
on the dense interior agree everywhere. Its source is
compact and its target is Hausdorff, so it is a
homeomorphism onto its image. Consequently (C.71.5)
has a unique continuous lift
$\overline R_{r,z}\in P^3$ through (C.71.7).
Continuity here includes the inverse maps, since the
source and target carry their homeomorphism-group
topologies. All these lifted maps fix the original
circle pointwise and centralize $Q$ on the interior.
They give the claimed based contraction. This argument
uses compact word control to prove boundary continuity;
it does not identify the full relative group topology
with the topology of an arbitrary subgroup of $N$.

\emph{The cyclic action and its exact disk coordinate.}
Write
\[
 \kappa(x)=\frac{x}{1+|x|},\qquad
 \kappa^{-1}(s)=\frac{s}{1-|s|}\quad(-1<s<1),
 \qquad S=[-1,1]^2.
                                                               \tag{C.71.8}
\]
The extension of $\kappa$ takes $-\infty,+\infty$ to
$-1,+1$. An explicit homeomorphism $J:S\to
\overline{\mathbb D}$ and its inverse are
\[
 \begin{aligned}
 J(v)&=\frac{\|v\|_\infty}{\|v\|_2}v\quad(v\ne0),
                          &J(0)&=0,\\
 J^{-1}(w)&=\frac{\|w\|_2}{\|w\|_\infty}w\quad(w\ne0),
                          &J^{-1}(0)&=0 .
 \end{aligned}                                                   \tag{C.71.9}
\]
The formulas are inverse on each ray and continuous at
zero. They take the square boundary onto the unit circle.
Retain the actual interior coordinate
$p=J\circ(\kappa\times\kappa):\mathbb R^2\to\mathbb D$.

On the real plane take $T(x,y)=(x,y+1)$.
Its extension on $S$ is
\[
 \overline T(s,t)=(s,\tau(t)),\qquad
 \tau(t)=\kappa(\kappa^{-1}(t)+1)\quad(|t|<1),
       \qquad\tau(-1)=-1,\quad\tau(1)=1.
                                                               \tag{C.71.10}
\]
The specified closed-disk generator is
$J\overline T J^{-1}$. The group
$Q=p\langle T\rangle p^{-1}$ is free and proper on
the interior: a compact real-plane set has bounded
$y$-coordinate and meets only finitely many of its
integer translates. It is closed and discrete in the
actual interior homeomorphism group, since evaluation
at $(0,0)$ sends $T^m$ to $(0,m)$; a convergent sequence
of such maps has eventually constant integer parameter.
Conjugation by $p$ transports these assertions, with
its inverse giving the opposite comparison.

The fixed set of $\overline T$ on $\partial S$ is
exactly the union of the two closed horizontal edges.
Indeed $\tau(t)>t$ for $-1<t<1$. Under a purported
closed-disk conjugacy $k$ to a generator $\gamma$ of
a torsion-free cyclic Fuchsian action, the equation
$k(J\overline T J^{-1})=\gamma k$ would imply
\[
       kJ\bigl(\operatorname{Fix}(\overline T)
                         \cap\partial S\bigr)
                  =\operatorname{Fix}(\gamma)\cap S^1.
                                                               \tag{C.71.11}
\]
The left side consists of two nondegenerate arcs,
whereas the right side has one point for a parabolic
element or two points for a hyperbolic element.
Since $kJ$ is injective, (C.71.11) is impossible.

There is a direct failure of the analogous partial
restriction comparison as well. Set
\[
 S_*=[-1,1]\times(-1,1),\qquad
 \Omega_*=\{-1,1\}\times(-1,1).
                                                               \tag{C.71.12}
\]
The $\langle\overline T\rangle$ action on $S_*$ is
free and proper. Through $(s,y)\mapsto(s,\kappa(y))$,
its actual quotient is the closed annulus
$[-1,1]\times(\mathbb R/\mathbb Z)$, with deck map
$y\mapsto y+1$ and boundary $\Omega_*/\langle\overline
T\rangle$. The horizontal translation $L(x,y)=(x+1,y)$
extends on $S$ as
$\overline L(s,t)=(\tau(s),t)$. It commutes with
$\overline T$, and its restriction to $S_*$ fixes
$\Omega_*$ pointwise. It does not fix the horizontal
edges pointwise. Thus
\[
 C_{\operatorname{Homeo}(S,\partial S)}(\overline T)
       \longrightarrow
 C_{\operatorname{Homeo}(S_*,\Omega_*)}(\overline T)
                                                               \tag{C.71.13}
\]
is not surjective: any preimage of $\overline L|_{S_*}$
would equal $\overline L$ on the dense set $S_*$ and
hence on $S$. The maps $J$ and $J^{-1}$ transport this
same counterexample to the specified disk.

\emph{The actual boundary-fixed centralizer of this action.}
Let $\mathcal V$ consist of continuous functions
$F:\mathbb R^2\to\mathbb R$, written $F_x(y)$, such that
each $F_x$ is an increasing homeomorphism of $\mathbb R$,
and
\[
 F_x(y+1)=F_x(y)+1,\qquad
 \lim_{|x|\to\infty}\sup_{0\leq y\leq1}
                         |F_x(y)-y|=0 .
                                                               \tag{C.71.14}
\]
Its product and inverse are the pointwise fiber
composition and inverse. Give it the metric topology
\[
 d_{\mathcal V}(F,G)=
  \sup_{x,y}|F_x(y)-G_x(y)|
   +\sup_{x,y}|F_x^{-1}(y)-G_x^{-1}(y)|.
                                                               \tag{C.71.15}
\]
Both suprema are finite. The displacements are
continuous and periodic in $y$, and vanish as
$|x|\to\infty$, so they are bounded globally.
The inverse displacements have the same suprema
at each $x$ and the same vanishing property.
Joint continuity of $(x,y)\mapsto F_x^{-1}(y)$
follows directly: on a compact $x$-interval the
periodic displacements are uniformly bounded,
which confines all inverse arguments to a compact
interval; strict monotonicity identifies every
subsequential limit.

There is an exact group isomorphism
\[
 \begin{split}
 \mathcal V&\longrightarrow
       C_{\operatorname{Homeo}(S,\partial S)}(\overline T),\\
 F&\longmapsto H_F,\qquad
 H_F(\kappa(x),\kappa(y))
          =(\kappa(x),\kappa(F_x(y))),
 \end{split}                                                   \tag{C.71.16}
\]
with the identity on $\partial S$. To prove surjectivity,
write the real-plane restriction of a map on the right
as $(X(x,y),Y(x,y))$. Commutation gives
$X(x,y+m)=X(x,y)$ and
$Y(x,y+m)=Y(x,y)+m$ for each integer $m$.
Letting $m\to+\infty$ and using the actual fixed top
edge forces $X(x,y)=x$. Bijectivity now gives a
homeomorphism $F_x(y)=Y(x,y)$ of each fiber, and the
integer equation forces it to be increasing.
Continuity at the two fixed vertical edges,
uniformly for $y\in[0,1]$, is exactly the vanishing
condition in (C.71.14).

Conversely (C.71.14) defines a homeomorphism of
the real plane with the stated inverse. For
bounded $x$ its displacement is uniformly bounded,
so it extends as the identity at the horizontal
edges. For $|x|\to\infty$ its displacement tends
uniformly to zero, giving the identity at the
vertical edges and at the corners. These assertions
also apply to the inverse. They prove the claimed
extension and bijection.

The isomorphism (C.71.16) is topological. In the
forward direction $\kappa$ is globally Lipschitz,
so each supremum in (C.71.15) controls the uniform
square displacement of the corresponding maps or
inverses. Conversely suppose $H_{F_j}$ and their
inverses converge uniformly to $H_F$ and its inverse.
For real $y\in[0,1]$, the values $F_x(y)$ lie in
one fixed compact real interval, uniformly in $x$.
Uniform convergence of their $\kappa$-coordinates
keeps all $F_{j,x}(y)$ in a slightly larger compact
real interval. Uniform continuity of $\kappa^{-1}$
there gives uniform convergence of these actual
values. Periodicity extends it to every real $y$.
Applying the same argument to the inverse maps
gives exactly (C.71.15).

Finally define
\[
       \mathcal H_r(F)_x(y)
                      =(1-r)F_x(y)+ry,\qquad 0\leq r\leq1 .
                                                               \tag{C.71.17}
\]
It is strictly increasing, proper, and has the
same integer equivariance and vanishing condition.
Forward continuity in (C.71.15) is immediate.
For inverse continuity, each fixed element of
$\mathcal V$ and its inverse is uniformly continuous
in the fiber variable, uniformly in $x$: its
displacement is continuous on
$\mathbb R\times(\mathbb R/\mathbb Z)$ and vanishes
at the two $x$-ends. The elementary increasing-map
inverse estimate therefore turns uniform forward
convergence to a member of $\mathcal V$ into uniform
inverse convergence. This applies to (C.71.17)
when $(r,F)$ varies. Hence it is a continuous
contraction of the full topological group to the
identity, fixing the identity throughout.
Conjugation by $J$ gives the stated disk contraction.

Every such map preserves the real $x$-coordinate.
For each specified $x$ and $c\in\mathbb R$, a
continuous compactly supported function $b$ with
$b(x)=c$ gives $F_u(y)=y+b(u)\in\mathcal V$.
Thus the exact group orbits in the interior are
the vertical lines. This records the transitivity
scope of the example explicitly.

The exact certificate
\path{code/part_b/c71_disk_compactification_certificate.py}
checks the rational compactification and inverse,
the translation conjugacies and commutation,
the displayed boundary and quotient maps, and
nonlinear periodic fiber models together with
their interpolation and inverse formulas.
The full topological comparisons and contractions
are proved above. This claim leaves the general
zero- and one-dimensional proper-operator disk
comparison unresolved.
\end{proof}
